% Generated draft: requires mathematical and source review.
\documentclass[11pt]{article}
\usepackage[margin=25mm]{geometry}
\usepackage{fontspec}
\setmainfont{Latin Modern Roman}
\usepackage{amsmath,amssymb,mathtools,unicode-math}
\setmathfont{Latin Modern Math}
\usepackage{enumitem,needspace,xurl,hyperref}
\setlength{\parindent}{0pt}
\setlength{\parskip}{5pt}
\setlength{\emergencystretch}{3em}
\newcommand{\N}{\mathbb{N}}
\newcommand{\Z}{\mathbb{Z}}
\newcommand{\Q}{\mathbb{Q}}
\newcommand{\R}{\mathbb{R}}
\newcommand{\C}{\mathbb{C}}
\newcommand{\F}{\mathbb{F}}
\newcommand{\A}{\mathbb{A}}
\newcommand{\PP}{\mathbb{P}}
\newcommand{\E}{\mathbb{E}}
\newcommand{\eps}{\varepsilon}
\newcommand{\dd}{\,\mathrm{d}}
\begin{document}
\section*{UnsolvedMath ID: 30003842}
\textbf{Source problem code:} catalogue entry
\subsection{Definitions and mathematical statement}
The supplied record does not determine a precise mathematical proposition. Its cleaned formulation asks whether Thompson's group \(F\) is the only ``structurally significant'' subgroup of \(\operatorname{PL}_o(I)\), described there as the group of orientation-preserving piecewise-linear homeomorphisms of an interval. Its original formulation instead asks whether \(F\) is the only ``interesting'' subgroup of a group described as consisting of piecewise-projective homeomorphisms, using the same notation. The supplied title also refers to piecewise-projective homeomorphisms. These descriptions do not specify the same ambient group.

For reference, a standard concrete realization of Thompson's group is
\[
F=\left\{f\colon[0,1]\longrightarrow[0,1]\;\middle|\;
\begin{array}{l}
f\text{ is an increasing piecewise-linear homeomorphism},\\
f\text{ has finitely many breakpoints, all in }\mathbb Z[1/2],\\
\text{every linear piece has slope in }\{2^k:k\in\mathbb Z\}
\end{array}\right\},
\]
with group operation composition, where
\[
\mathbb Z[1/2]=\{m/2^n:m\in\mathbb Z,\ n\in\mathbb Z_{\geq 0}\}.
\]
In this realization, every element fixes \(0\) and \(1\). Under the convention of finitely many linear pieces, the group of all increasing piecewise-linear homeomorphisms of \([0,1]\) contains this realization of \(F\). This supplies a standard reference model; the record does not explicitly specify that \(I=[0,1]\), its endpoint conditions, or its finiteness convention. Piecewise-projective maps instead have pieces given by fractional linear transformations, so the original wording cannot silently be replaced by the piecewise-linear interpretation.

\textbf{Unresolved statement.} Determine the intended ambient group \(G\), the meaning of ``interesting'' or ``structurally significant'' for a subgroup \(H\leq G\), and the intended meaning of uniqueness. In particular, the record does not specify whether uniqueness means equality with a distinguished embedded copy of \(F\), conjugacy to such a copy, or isomorphism to \(F\). Without these specifications, no definite quantified uniqueness assertion can be extracted. The supplied label \texttt{open} describes the record's curation status; it does not establish a well-defined open conjecture.
\subsection{Short English statement}
The record asks whether Thompson's group \(F\) is the only subgroup of a certain homeomorphism group deserving an unspecified structural distinction. A precise conjecture requires clarification of that distinction, the ambient piecewise-linear versus piecewise-projective group, the interval conventions, and the meaning of uniqueness.
\subsection{Sources}
\begin{itemize}
\item Supplied local UnsolvedMath registry record, permanent ID \(30003842\), source problem code \texttt{catalogue entry}, collection \emph{Oberwolfach Reports Open Problems}. No individual proposer is supplied; the recorded proposed year is \(2018\).
\item Supplied report citation: \emph{Cohomological and Metric Properties of Groups of Homeomorphisms of R (2019)}, Oberwolfach Reports, pp.~1579--1633, DOI \texttt{10.4171/owr/2018/26}. \url{https://doi.org/10.4171/owr/2018/26}. The supplied background describes these reports as documenting mathematical workshops and their open-problem discussions.
\item Supplied literature source labeled ``Brin'': \url{https://arxiv.org/abs/math/0501085}. No further bibliographic details or supporting passage are supplied.
\item Supplied literature source labeled ``Oberwolfach Report 2019'': \url{https://ems.press/content/serial-article-files/46748}.
\item Literature assessment dated exactly \texttt{2026-08-21}: the supplied curation labels the problem open, observes that ``interesting subgroup'' is not a precise theorem, and describes the subgroup structure of \(\operatorname{PL}_o(I)\) as poorly understood in the primary literature. This is supplied dated curation derived from the OWR material, not an independently verified assessment. The supplied background expressly requests independent verification of time-sensitive status and cited claims.
\item The record labels its cleaned statement \texttt{corrected\_verified} and reports a manual semantic, typographic, and LaTeX review asserting preservation of the mathematical target. This is supplied review metadata; it does not resolve the formulation conflicts identified here. No authoritative source passage resolving those conflicts is supplied.
\end{itemize}
\end{document}
